\documentclass[11pt]{article}
\usepackage[utf8]{inputenc}
\usepackage[T1]{fontenc}
\usepackage{amsmath,amsfonts,amssymb}
\usepackage{geometry}
\usepackage{natbib}
\usepackage{url}
\usepackage{graphicx}
\usepackage[breaklinks=true]{hyperref}
\usepackage{booktabs}
\usepackage{longtable}
\usepackage{array}

\geometry{margin=1in}
\setlength{\tabcolsep}{4pt}
\makeatletter\g@addto@macro{\UrlBreaks}{\do\_}\makeatother
\emergencystretch=1.5em

\newcommand{\degC}{\ensuremath{^{\circ}}C}
\newcommand{\Kop}{K\"oppen}
% Figure placeholder: shows the figure if the file exists, otherwise a labelled box.
\newcommand{\figholder}[3]{%
  \begin{figure}[htbp]\centering
  \IfFileExists{#1}{\includegraphics[width=\textwidth,height=0.42\textheight,keepaspectratio]{#1}}{%
    \fbox{\parbox{0.92\textwidth}{\centering\vspace{1.6cm}\small Figure placeholder: \texttt{\detokenize{#1}}\\
    generated by \texttt{\detokenize{code/make_figures.py}} (or the pipeline, where noted)\vspace{1.6cm}}}}
  \caption{#2}\label{#3}
  \end{figure}}

\title{\textbf{How Much Warmer? Plant Range Extremes as Bounds on Holocene Temperature in the Lowlands of Washington State}}

\author{David D. Zelenka}

\date{\today\\v.0.3 (data draft)}

\begin{document}

\maketitle

\begin{abstract}
Glaciers record the cold extremes of the Holocene in the Pacific Northwest, most clearly the Little Ice Age, but the warm extremes leave little physical trace outside lake sediments and the distributions of living plants. This paper asks how much warmer than the present the warm periods of the Holocene can have been in the lowlands of Washington State, and uses plant range extremes to answer it. The logic is a bound, not a reconstruction: a taxon that persisted at a site through the Holocene cannot have been pushed beyond the climatic edge of its range, and a taxon that grows nearby but never colonised the site cannot have had its climate there. Range edges were computed for 26 indicator taxa from latitude-stratified, cleaned GBIF occurrences and WorldClim~2.1 temperature (1970--2000), and combined with the Holocene status of taxa at eight sediment records. In the western lowlands, persistent Douglas-fir caps Holocene summer warming at about $+2.5$\,\degC{} relative to 1970--2000 at Lake Washington and $+4.1$ to $+5.0$\,\degC{} at four other sites; these caps are provisional because Douglas-fir is limited mainly by moisture. On the cold side, the absence of mountain hemlock from Holocene lowland records caps summer cooling at about $-2.5$\,\degC{}, a bound set by a temperature-limited taxon. No taxon was lost or gained outright at any site, so the data give no hard lower bound on warming: the known warm periods lie between the present and the cap. Garry oak and Pacific madrone add a sharp local constraint on the cold side, because the north Olympic and San Juan rainshadow lies at the cool-summer edge of both species. As a test of the underlying idea, eastern lowland taxa have broader summer-temperature ranges than western lowland or alpine taxa (median difference 3.7\,\degC, permutation $p=0.002$). The sample is small and several bounds are provisional; the methods, data, and code are offered as a starting point for further work.
\end{abstract}

\section{Introduction}\label{sec:intro}

The cold side of Holocene climate in the Pacific Northwest is well recorded. Glacier moraines show Neoglacial readvances culminating in the Little Ice Age (LIA), and lake cores in the Olympic Mountains record the cooling and moistening that followed the early Holocene \citep{menounos2009,gavin2001,zelenka2026olympic}. The warm side is harder to measure. Warm intervals erase rather than build glacial landforms, so their evidence is indirect: shifts in pollen and charcoal in lake sediments \citep{whitlock1992,walsh2008battleground,leopold2016orcas}, and the present distributions of plants that may have spread during warmer times.

This paper focuses on that warm side, with one question: \emph{how much warmer than today can the Holocene warm periods have been in Washington's lowlands?} A lower bound, how much warmer they \emph{must} have been, is reported where the data allow.

The approach rests on a simple proposition. If Holocene climate had never varied, the plants of a site would need to tolerate only one climate; if it varied widely, the plants that remain must tolerate all of it. Over six to ten thousand years and many generations, a taxon intolerant of the warm or cold extremes would be lost. The taxa that persisted therefore carry, in their range extremes, a ceiling on how far the climate can have moved. The idea is related to the climate variability hypothesis \citep{janzen1967,stevens1989,ghalambor2006} and to the oscillation hypothesis \citep{dynesiusjansson2000}, and its logic follows the indicator-species and mutual-climatic-range traditions \citep{iversen1944,atkinson1987,thompson2012atlasF}.

Temperature is treated alone in this draft. Precipitation data were not available for the analysis, and the forcing of Holocene variability (orbital insolation, ocean circulation, or both) is left open \citep{bartlein1998,barron2003}.

\section{The Logic of Bounds}\label{sec:logic}

Let $x_0$ be a site's modern value of a climate variable and $x(t)$ its value at time $t$ in the Holocene. For a taxon $s$, let $L_s$ and $U_s$ be the lower and upper edges of its range in that variable, taken as quantiles of the variable across the taxon's occurrences. Each taxon recorded at a site has one of four Holocene statuses, and each constrains $x(t)$ differently.
\begin{description}
\item[Persistent] (present throughout). Every Holocene climate lay within the range, so $L_s \le x(t) \le U_s$. This caps warming at $U_s - x_0$ and cooling at $x_0 - L_s$.
\item[Never present] (grows within dispersal distance today, never recorded at the site). The climate never entered its range. A cold-adapted taxon with $U_s < x_0$ caps cooling at $x_0 - U_s$; a warm-adapted taxon with $L_s > x_0$ caps warming at $L_s - x_0$.
\item[Lost] (present earlier, absent now). The climate must once have been within the range, giving a lower bound on the excursion.
\item[Gained] (absent earlier, present now). Unless the arrival reflects slow migration, the climate must once have been outside the range, again a lower bound.
\end{description}
For a site with persistent taxa $P$ and never-present taxa $N$, the upper bound on warming is
\begin{equation}
\Delta^{+}_{\max} = \min\Big( \min_{s\in P} U_s,\; \min_{s\in N,\,L_s>x_0} L_s \Big) - x_0 ,
\label{eq:upmax}
\end{equation}
and the bound on cooling is symmetric. The quantile used for each edge is the one that keeps the bound conservative: the 1st and 99th percentiles for persistent, lost, and gained taxa, and the 5th and 95th for never-present taxa. Bounds are computed twice, from the whole species range and from records within 800~km of the site, because a species-wide range is wider than the tolerance of local populations.

Each bound is graded. It is \emph{ok} when the taxon that sets it is thought to be limited by that variable; \emph{provisional} when that taxon is limited mainly by something else, usually moisture; and \emph{not informative} (GEO) when the taxon's edge coincides with the edge of the land climate that exists within 800~km. A coastal tree whose warmest-winter edge is the coastline, for example, says nothing about winter tolerance.

For a statewide picture, a spatially uniform change is assumed, so that the smallest cap among sites binds:
\begin{equation}
\Delta^{+}_{\text{region}} = \min_{k} \Delta^{+}_{\max,k}.
\label{eq:region}
\end{equation}

\section{Method Check}\label{sec:check}

Table~\ref{tab:check} lists the problems identified in the method and how each is handled. Two deserve emphasis for the warm-side question.

First, \emph{the warm caps rest on moisture-limited trees}. Douglas-fir, western hemlock, and Sitka spruce are absent from the hot interior lowlands mainly because summers there are dry, not because they are hot. Their warm-temperature edges therefore partly record a moisture edge, and warm caps set by them are graded provisional. Adding precipitation, or finding warm-side caps set by temperature-limited taxa, is the main step needed to firm them up (Section~\ref{sec:next}).

Second, \emph{longevity}. A Garry oak can live for about five centuries, so a single cold or warm century need not remove a population. The bounds therefore apply to climate sustained over many generations, not to individual decades. Over the six to ten thousand years these records span, with repeated warm and cold episodes, persistence through every episode is strong evidence that no sustained excursion exceeded the taxon's range.

{\small
\begin{longtable}{>{\raggedright\arraybackslash}p{0.05\textwidth}>{\raggedright\arraybackslash}p{0.25\textwidth}>{\raggedright\arraybackslash}p{0.29\textwidth}>{\raggedright\arraybackslash}p{0.33\textwidth}}
\caption{Method check: issues and treatment.}\label{tab:check}\\
\toprule
\textbf{ID} & \textbf{Issue} & \textbf{Effect if ignored} & \textbf{Treatment} \\
\midrule
\endfirsthead
\toprule
\textbf{ID} & \textbf{Issue} & \textbf{Effect if ignored} & \textbf{Treatment} \\
\midrule
\endhead
\bottomrule
\endfoot
C1 & Species are older than the Holocene & Species-wide breadth reflects glacial-cycle variability & Holocene signal taken from site status (persistence, absence), not evolution \citep{smithbeaulieu2009,ackerly2003} \\
C2 & Tracking vs.\ tolerance & Local migration mimics persistence & Lowland focus (gentle terrain); local records preferred \citep{davisshaw2001,rull2009} \\
C3 & Spatial vs.\ temporal breadth & Widespread species look tolerant & Local breadth; range-size residuals \citep{slatyer2013} \\
C4 & Realised vs.\ fundamental niche; ecotypes & Bounds too tight or too loose & Species-wide and 800~km envelopes; GEO check \citep{sax2013,rehfeldt1999} \\
C5 & Soils, fire, people, dispersal & Spurious bounds & Status records graded; high-risk records excluded \\
C6 & Taxonomy, pollen lumping & Wrong envelope & Explicit fossil-to-species mapping; \emph{Picea} and Ericaceae limitations stated \\
C7 & Abundance vs.\ presence & Overstated bounds & Abundance changes excluded from bounds \\
C8 & Temperature confounded with moisture & Moisture edge read as temperature edge & Bounds graded ok / provisional \\
C9 & Geographic truncation & Coastline read as tolerance & GEO grade \\
C10 & Longevity & Short excursions undetected & Bounds stated for sustained, multi-generation climate \\
C11 & Planted and misplaced records & False range edges & Isolation filter (5th neighbour $>$200~km): 373 records removed \\
C12 & Grid resolution and microsites & Site values too cold or warm & 2.5$'$ grid; 30$''$ rerun recommended; Hoh and Wentworth coordinates to be verified \\
\end{longtable}
}

\section{Data and Methods}\label{sec:data}

\paragraph{Occurrences.} For each taxon, GBIF records were sampled across North America (15--72$^\circ$N, 170--52$^\circ$W) in 1$^\circ$ latitude bands, up to 150 records per band at random positions, so that range edges are sampled regardless of how densely the core is recorded. Records with coordinate uncertainty above 5~km, introduced or managed records, and duplicates were removed, followed by spatially isolated records (Table~\ref{tab:taxa}).

\paragraph{Climate.} Monthly mean temperature from WorldClim~2.1 at 2.5$'$ (about 4.5~km), 1970--2000 \citep{fick2017}, gives mean annual temperature (MAT), coldest- and warmest-month means (MTCO, MTWA), and growing degree-days above 5\,\degC{} (GDD5). All departures in this paper are relative to 1970--2000. The bounds use temperature only; monthly precipitation from the same source is used solely to classify climates on the maps, where it is held at its present values.

\paragraph{Sites and status.} The Holocene status of taxa at each site was compiled from the sediment literature (Table~\ref{tab:status}). Site climates are in Table~\ref{tab:sites}.

\paragraph{Code.} All results come from \texttt{build\_variability\_proxy.py}; all tables in this paper are generated from its output by \texttt{make\_tables.py}, and all figures by \texttt{make\_figures.py}.

{\small
\begin{longtable}{>{\raggedright\arraybackslash}p{0.21\textwidth}>{\raggedright\arraybackslash}p{0.10\textwidth}rr>{\raggedright\arraybackslash}p{0.09\textwidth}>{\raggedright\arraybackslash}p{0.10\textwidth}>{\raggedright\arraybackslash}p{0.16\textwidth}}
\caption{Indicator taxa and occurrence data. n: cleaned, climate-annotated records after latitude-stratified sampling; Drop.: spatially isolated records removed (likely planted or mis-georeferenced). Last column: variables the taxon is thought to be limited by (used to grade bounds as \emph{ok} or \emph{provisional}).}\label{tab:taxa}\\
\toprule
\textbf{Taxon} & \textbf{Group} & \textbf{n} & \textbf{Drop.} & \textbf{Lat.\ 5--95\%} & \textbf{MTWA 5--95\% (\degC)} & \textbf{Thought limited by} \\
\midrule
\endfirsthead
\toprule
\textbf{Taxon} & \textbf{Group} & \textbf{n} & \textbf{Drop.} & \textbf{Lat.\ 5--95\%} & \textbf{MTWA 5--95\% (\degC)} & \textbf{Thought limited by} \\
\midrule
\endhead
\bottomrule
\endfoot
\emph{Quercus garryana} & west lowland & 1719 & 1 & 35.5--49.1 & 16.4--23.8 & MTWA; GDD5; P\_JJA \\
\emph{Arbutus menziesii} & west lowland & 1920 & 1 & 34.5--50.0 & 16.0--22.3 & MTCO; P\_JJA \\
\emph{Arctostaphylos columbiana} & west lowland & 1089 & 1 & 38.8--50.0 & 14.7--18.7 & P\_JJA \\
\emph{Juniperus maritima} & west lowland & 684 & 0 & 48.4--49.6 & 15.3--17.5 & MI \\
\emph{Opuntia fragilis} & west lowland & 1550 & 16 & 36.7--55.3 & 16.1--23.8 & MI (not T) \\
\emph{Pseudotsuga menziesii} & west lowland & 2865 & 28 & 24.1--53.1 & 13.6--22.0 & MI; P\_JJA \\
\emph{Thuja plicata} & west lowland & 1735 & 28 & 40.9--55.3 & 13.3--19.2 & MI; P\_JJA \\
\emph{Tsuga heterophylla} & west lowland & 2201 & 2 & 39.4--59.3 & 12.2--18.2 & MI; P\_JJA \\
\emph{Picea sitchensis} & west lowland & 2058 & 20 & 40.9--59.6 & 12.2--17.4 & P\_JJA; MTCO \\
\emph{Rhododendron macrophyllum} & west lowland & 1109 & 3 & 38.6--49.1 & 14.1--18.2 & MTCO; MTWA (weak) \\
\emph{Rhododendron groenlandicum} & west lowland & 2539 & 59 & 43.5--65.3 & 11.0--19.8 & MTWA (weak) \\
\emph{Artemisia tridentata} & east lowland & 2317 & 2 & 31.8--50.7 & 14.4--24.9 & MI; MAP \\
\emph{Purshia tridentata} & east lowland & 1763 & 9 & 35.9--49.0 & 14.7--23.6 & MI \\
\emph{Sarcobatus vermiculatus} & east lowland & 1469 & 6 & 35.1--50.3 & 17.5--26.9 & MI \\
\emph{Grayia spinosa} & east lowland & 1264 & 2 & 34.0--47.0 & 17.5--27.8 & MI; MTWA \\
\emph{Celtis reticulata} & east lowland & 1554 & 18 & 23.9--45.6 & 20.6--29.5 & MTWA; GDD5 \\
\emph{Juniperus occidentalis} & east lowland & 1164 & 10 & 34.3--46.4 & 13.4--23.5 & MI \\
\emph{Pinus ponderosa} & east lowland & 2109 & 37 & 33.0--51.0 & 14.7--22.2 & MI; MAP \\
\emph{Opuntia polyacantha} & east lowland & 2118 & 27 & 32.1--50.8 & 16.9--27.3 & MI \\
\emph{Tsuga mertensiana} & alpine ref. & 2033 & 1 & 37.7--60.9 & 10.3--15.3 & MTWA; GDD5 \\
\emph{Abies lasiocarpa} & alpine ref. & 2565 & 10 & 37.4--60.1 & 9.3--16.2 & MTWA; GDD5 \\
\emph{Abies amabilis} & alpine ref. & 1043 & 6 & 43.6--54.2 & 11.6--16.2 & MTWA \\
\emph{Callitropsis nootkatensis} & alpine ref. & 1189 & 3 & 44.1--57.1 & 11.7--15.8 & MTWA \\
\emph{Pinus albicaulis} & alpine ref. & 1565 & 9 & 37.5--53.0 & 9.3--15.0 & MTWA; GDD5 \\
\emph{Picea engelmannii} & alpine ref. & 1759 & 12 & 31.9--51.8 & 9.3--17.9 & MTWA \\
\emph{Alnus rubra} & disturbance & 1876 & 19 & 37.9--57.1 & 13.0--18.7 & none \\
\emph{Pteridium aquilinum} & disturbance & 4113 & 45 & 18.8--54.1 & 14.0--28.0 & none \\
\end{longtable}
}


\figholder{figures/fig_sites.png}{Sediment records used. Marker size: number of taxon status records at the site. Dashed: approximate Cascade crest.}{fig:sites}

{\small
\begin{longtable}{>{\raggedright\arraybackslash}p{0.22\textwidth}crrrrrr}
\caption{Modern (1970--2000) climate at the sediment sites, WorldClim 2.1 at 2.5$'$. Temperatures in \degC. Coordinates for Hoh and Wentworth are approximate and place those sites too high and cold.}\label{tab:sites}\\
\toprule
\textbf{Site} & \textbf{Side} & \textbf{Lat} & \textbf{Lon} & \textbf{MAT} & \textbf{MTCO} & \textbf{MTWA} & \textbf{GDD5} \\
\midrule
\endfirsthead
\toprule
\textbf{Site} & \textbf{Side} & \textbf{Lat} & \textbf{Lon} & \textbf{MAT} & \textbf{MTCO} & \textbf{MTWA} & \textbf{GDD5} \\
\midrule
\endhead
\bottomrule
\endfoot
Battle Ground Lake & W & 45.80 & -122.49 & 10.5 & 3.5 & 18.1 & 2096 \\
Lake Washington & W & 47.62 & -122.26 & 11.1 & 5.0 & 18.4 & 2224 \\
Killebrew Lake Fen & W & 48.60 & -122.90 & 9.5 & 4.2 & 15.8 & 1709 \\
Pender Is.\ (BC) & W & 48.80 & -123.30 & 9.7 & 4.3 & 16.2 & 1776 \\
Hoh valley & W & 47.86 & -124.00 & 7.7 & 2.2 & 14.4 & 1238 \\
Crocker Lake & W & 47.94 & -122.88 & 9.5 & 3.6 & 16.4 & 1742 \\
Wentworth Lake & W & 48.00 & -124.50 & 9.5 & 4.3 & 15.6 & 1697 \\
Williams Lake Fen & E & 47.32 & -117.62 & 8.7 & -1.6 & 20.0 & 1988 \\
Wildcat Lake & E & 46.90 & -117.40 & 9.0 & -0.7 & 19.7 & 1982 \\
Carp Lake & E & 45.92 & -120.88 & 7.3 & -2.0 & 17.6 & 1547 \\
Castor Lake & E & 48.54 & -119.56 & 8.7 & -4.3 & 21.1 & 2187 \\
\end{longtable}
}


{\small
\begin{longtable}{>{\raggedright\arraybackslash}p{0.20\textwidth}>{\raggedright\arraybackslash}p{0.25\textwidth}>{\raggedright\arraybackslash}p{0.16\textwidth}>{\raggedright\arraybackslash}p{0.14\textwidth}l}
\caption{Holocene status of taxa at the sites, compiled from the sediment literature (sources in \texttt{taxon\_site\_status.csv}). Only \emph{persistent} and \emph{never present} enter upper bounds; \emph{lost} and \emph{gained} would give lower bounds; \emph{abundance change} gives no hard bound.}\label{tab:status}\\
\toprule
\textbf{Site} & \textbf{Taxon} & \textbf{Holocene status} & \textbf{Confidence} & \textbf{In bounds} \\
\midrule
\endfirsthead
\toprule
\textbf{Site} & \textbf{Taxon} & \textbf{Holocene status} & \textbf{Confidence} & \textbf{In bounds} \\
\midrule
\endhead
\bottomrule
\endfoot
Battle Ground Lake & \emph{Pseudotsuga menziesii} & persistent & medium (regional pollen) & used \\
Battle Ground Lake & \emph{Quercus garryana} & abundance change & high & used \\
Battle Ground Lake & \emph{Thuja plicata} & abundance change & medium & used \\
Battle Ground Lake & \emph{Tsuga heterophylla} & abundance change & medium & used \\
Battle Ground Lake & \emph{Tsuga mertensiana} & never present & medium & used \\
Lake Washington & \emph{Pseudotsuga menziesii} & persistent & medium & used \\
Lake Washington & \emph{Thuja plicata} & abundance change & medium & used \\
Lake Washington & \emph{Tsuga mertensiana} & never present & medium & used \\
Lake Washington & \emph{Picea engelmannii} & never present & low & excluded \\
Killebrew Lake Fen & \emph{Pseudotsuga menziesii} & persistent & low-medium & used \\
Killebrew Lake Fen & \emph{Quercus garryana} & abundance change & high & used \\
Hoh valley & \emph{Tsuga heterophylla} & persistent & medium & used \\
Hoh valley & \emph{Picea sitchensis} & persistent & medium & used \\
Hoh valley & \emph{Pseudotsuga menziesii} & abundance change & medium & used \\
Pender Is.\ (BC) & \emph{Pseudotsuga menziesii} & persistent & medium & used \\
Pender Is.\ (BC) & \emph{Quercus garryana} & abundance change & high & used \\
Carp Lake & \emph{Artemisia tridentata} & abundance change & medium & used \\
Carp Lake & \emph{Pinus ponderosa} & abundance change & medium & used \\
Carp Lake & \emph{Quercus garryana} & gained & low-medium & excluded \\
Williams Lake Fen & \emph{Artemisia tridentata} & persistent & medium & used \\
Williams Lake Fen & \emph{Pinus ponderosa} & abundance change & medium & used \\
Wildcat Lake & \emph{Artemisia tridentata} & persistent & medium & used \\
Wildcat Lake & \emph{Pinus ponderosa} & never present & low & used \\
\end{longtable}
}


\section{Results}\label{sec:results}

\subsection{How much warmer: caps at each site}

Table~\ref{tab:bounds} and Figure~\ref{fig:bounds} give the caps at each site. In the western lowlands every upward cap on summer temperature is set by a persistent conifer, Douglas-fir at four sites and Sitka spruce at the Hoh, and every one is provisional. The regional caps on summer warming are $+2.5$\,\degC{} at Lake Washington and $+4.1$ to $+5.0$\,\degC{} at the other four western sites; species-wide caps are wider, $+4.0$ to $+8.7$\,\degC. Most upward caps on mean annual temperature are graded GEO, because the persistent coastal conifers already occupy the mildest climates available near the coast. In the east, the only persistent taxon is sagebrush, whose regional warm caps are GEO and whose species-wide caps are provisional.

{\footnotesize
\begin{longtable}{>{\raggedright\arraybackslash}p{0.12\textwidth}lrrrr>{\raggedright\arraybackslash}p{0.11\textwidth}>{\raggedright\arraybackslash}p{0.11\textwidth}}
\caption{Maximum Holocene departure from the 1970--2000 mean permitted by each site's taxa (\degC). \emph{regional}: envelopes from records within 800~km; \emph{species}: whole range. Grades: \textsuperscript{ok} set by a taxon thought to be limited by that variable; \textsuperscript{prov} provisional (limiting taxon limited mainly by something else); \textsuperscript{GEO} not informative (the taxon's edge is the edge of land climate available within 800~km). Carp Lake has no usable status records.}\label{tab:bounds}\\
\toprule
\textbf{Site} & \textbf{Envelope} & \textbf{MTWA up} & \textbf{MTWA down} & \textbf{MAT up} & \textbf{MAT down} & \textbf{Sets MTWA up} & \textbf{Sets MTWA down} \\
\midrule
\endfirsthead
\toprule
\textbf{Site} & \textbf{Envelope} & \textbf{MTWA up} & \textbf{MTWA down} & \textbf{MAT up} & \textbf{MAT down} & \textbf{Sets MTWA up} & \textbf{Sets MTWA down} \\
\midrule
\endhead
\bottomrule
\endfoot
Battle Ground Lake & regional & $+$4.9 \textsuperscript{prov} & $-$2.5 \textsuperscript{ok} & $+$3.2 \textsuperscript{prov} & $-$3.3 \textsuperscript{ok} & \emph{P.\ menziesii} & \emph{T.\ mertensiana} \\
 & species & $+$6.5 \textsuperscript{prov} & $-$2.7 \textsuperscript{ok} & $+$7.1 \textsuperscript{prov} & $-$3.5 \textsuperscript{ok} & \emph{P.\ menziesii} & \emph{T.\ mertensiana} \\
Lake Washington & regional & $+$2.5 \textsuperscript{prov} & $-$3.2 \textsuperscript{ok} & $+$1.0 \textsuperscript{GEO} & $-$3.8 \textsuperscript{ok} & \emph{P.\ menziesii} & \emph{T.\ mertensiana} \\
 & species & $+$6.1 \textsuperscript{prov} & $-$3.1 \textsuperscript{ok} & $+$6.6 \textsuperscript{prov} & $-$4.1 \textsuperscript{ok} & \emph{P.\ menziesii} & \emph{T.\ mertensiana} \\
Killebrew Lake Fen & regional & $+$5.0 \textsuperscript{prov} & $-$3.8 \textsuperscript{prov} & $+$2.5 \textsuperscript{GEO} & $-$8.4 \textsuperscript{prov} & \emph{P.\ menziesii} & \emph{P.\ menziesii} \\
 & species & $+$8.7 \textsuperscript{prov} & $-$3.8 \textsuperscript{prov} & $+$8.1 \textsuperscript{prov} & $-$8.4 \textsuperscript{prov} & \emph{P.\ menziesii} & \emph{P.\ menziesii} \\
Pender Is.\ (BC) & regional & $+$4.6 \textsuperscript{prov} & $-$4.2 \textsuperscript{prov} & $+$2.3 \textsuperscript{GEO} & $-$8.6 \textsuperscript{prov} & \emph{P.\ menziesii} & \emph{P.\ menziesii} \\
 & species & $+$8.3 \textsuperscript{prov} & $-$4.2 \textsuperscript{prov} & $+$7.9 \textsuperscript{prov} & $-$8.6 \textsuperscript{prov} & \emph{P.\ menziesii} & \emph{P.\ menziesii} \\
Hoh valley & regional & $+$4.1 \textsuperscript{prov} & $-$1.1 \textsuperscript{prov} & $+$4.3 \textsuperscript{GEO} & $-$2.0 \textsuperscript{ok} & \emph{P.\ sitchensis} & \emph{P.\ sitchensis} \\
 & species & $+$4.0 \textsuperscript{prov} & $-$3.5 \textsuperscript{prov} & $+$4.3 \textsuperscript{ok} & $-$6.7 \textsuperscript{ok} & \emph{P.\ sitchensis} & \emph{P.\ sitchensis} \\
Williams Lake Fen & regional & $+$3.6 \textsuperscript{GEO} & $-$8.3 \textsuperscript{prov} & $+$3.2 \textsuperscript{GEO} & $-$7.9 \textsuperscript{prov} & \emph{A.\ tridentata} & \emph{A.\ tridentata} \\
 & species & $+$6.1 \textsuperscript{prov} & $-$7.7 \textsuperscript{prov} & $+$8.1 \textsuperscript{prov} & $-$7.7 \textsuperscript{prov} & \emph{A.\ tridentata} & \emph{A.\ tridentata} \\
Wildcat Lake & regional & $+$4.2 \textsuperscript{GEO} & $-$8.1 \textsuperscript{prov} & $+$2.9 \textsuperscript{GEO} & $-$8.2 \textsuperscript{prov} & \emph{A.\ tridentata} & \emph{A.\ tridentata} \\
 & species & $+$6.4 \textsuperscript{prov} & $-$7.5 \textsuperscript{prov} & $+$7.8 \textsuperscript{prov} & $-$8.0 \textsuperscript{prov} & \emph{A.\ tridentata} & \emph{A.\ tridentata} \\
\end{longtable}
}


\figholder{figures/fig_bounds.png}{Caps on Holocene departure at each site, warmest-month mean (left) and mean annual temperature (right), relative to 1970--2000. Right of zero: how much warmer; left: how much colder. Colour gives the grade; thick bars use 800~km envelopes, thin bars whole-range envelopes. Dotted: $\pm2.5$\,\degC.}{fig:bounds}

\subsection{How much colder: the cold-side cap}

The best-graded result in the dataset is on the cold side. Mountain hemlock grows 50--100~km from Battle Ground Lake and Lake Washington but did not colonise either site during the Holocene. Because mountain hemlock is limited by summer warmth, its absence caps summer cooling at $-2.5$\,\degC{} at Battle Ground Lake and $-3.1$ to $-3.2$\,\degC{} at Lake Washington, and mean annual cooling at $-3.3$ to $-4.1$\,\degC. These caps apply to the whole Holocene, including the LIA, and agree between regional and species-wide envelopes.

\subsection{Regional scenarios}

Under a uniform change across the western lowlands, the tightest site binds (Equation~\ref{eq:region}; Table~\ref{tab:scenarios}). Holocene summers in the western lowlands were therefore no more than about 2.5\,\degC{} cooler (graded ok) and, provisionally, no more than about 2.5\,\degC{} warmer than in 1970--2000. The warm figure is set by one site, Lake Washington; at the four other western sites the cap is roughly twice as large.

{\footnotesize
\begin{longtable}{>{\raggedright\arraybackslash}p{0.11\textwidth}r>{\raggedright\arraybackslash}p{0.11\textwidth}l>{\raggedright\arraybackslash}p{0.12\textwidth}l>{\raggedright\arraybackslash}p{0.24\textwidth}}
\caption{Uniform-shift scenarios for the western lowlands. Under a spatially uniform change, the smallest cap among sites binds. Cold: smallest downward cap graded \emph{ok}. Warm: smallest upward cap not graded GEO.}\label{tab:scenarios}\\
\toprule
\textbf{Scenario} & \textbf{MTWA shift} & \textbf{Binding site} & \textbf{Env.} & \textbf{Binding taxon} & \textbf{Grade} & \textbf{All western caps (smallest per site)} \\
\midrule
\endfirsthead
\toprule
\textbf{Scenario} & \textbf{MTWA shift} & \textbf{Binding site} & \textbf{Env.} & \textbf{Binding taxon} & \textbf{Grade} & \textbf{All western caps (smallest per site)} \\
\midrule
\endhead
\bottomrule
\endfoot
Coldest permitted & $-$2.50 & Battle Ground Lake & regional & \emph{Tsuga mertensiana} & ok & Battle Ground Lake 2.5, Lake Washington 3.1 \\
Warmest permitted & $+$2.46 & Lake Washington & regional & \emph{Pseudotsuga menziesii} & provisional & Lake Washington 2.5, Hoh valley 4.0, Pender Is.\ (BC) 4.6, Battle Ground Lake 4.9, Killebrew Lake Fen 5.0 \\
\end{longtable}
}


\subsection{Is there a lower bound?}

No. None of the 23 status records describes a taxon lost or gained outright; every documented Holocene change at these sites is a change in abundance, such as the expansion of Garry oak savanna at Battle Ground Lake between about 10.8 and 5.2~cal~ka \citep{walsh2008battleground} and on Pender Island after the Mazama eruption \citep{lucaslacourse2013}. Such expansions are consistent with warmer or drier summers, but they set no minimum on the magnitude. The data therefore place the known Holocene warm periods between the present and the cap: warmer than 1970--2000 by an amount between 0 and roughly 2.5\,\degC{} in summer at the most constrained site, and between 0 and 4--5\,\degC{} at the others.

\subsection{Garry oak and madrone: the rainshadow at the cool edge}

Garry oak and Pacific madrone, the two Mediterranean-affinity trees of the region, sharpen the cold side locally (Table~\ref{tab:oakmadrone}, Figure~\ref{fig:oakmadrone}). The warmest-month means at Killebrew Lake, Pender Island, and Crocker Lake lie at or below the 5th percentile of Garry oak's range and within 0.3\,\degC{} of madrone's. If both have persisted in the rainshadow, as their present populations and the oak pollen records suggest, sustained summers there cannot have been much cooler than in 1970--2000. Madrone is also frost-sensitive: its coldest-month edge leaves only 1.3--2.1\,\degC{} of room for winter cooling, the only winter-side bound in the dataset set by a taxon limited by winter cold. Neither species constrains warming; their warm edges lie 6--8\,\degC{} higher in summer. Two caveats apply. Both trees occupy warm south-facing microsites that the 4.5~km grid averages away, so the true room is somewhat larger; and madrone pollen cannot be identified below Ericaceae, so its local persistence rests on living populations. Pacific rhododendron has been added to the taxon list but not yet analysed.

{\small
\begin{longtable}{>{\raggedright\arraybackslash}p{0.22\textwidth}lrrrr}
\caption{Room for cooling / room for warming (\degC; GDD5 in degree-days) before the north Olympic and San Juan rainshadow sites leave the 5--95\% range of each species. Negative room: the site already lies beyond the species' 5th percentile.}\label{tab:oakmadrone}\\
\toprule
\textbf{Taxon} & \textbf{Var.} & \textbf{5--95\% range} & \textbf{Killebrew} & \textbf{Pender} & \textbf{Crocker} \\
\midrule
\endfirsthead
\toprule
\textbf{Taxon} & \textbf{Var.} & \textbf{5--95\% range} & \textbf{Killebrew} & \textbf{Pender} & \textbf{Crocker} \\
\midrule
\endhead
\bottomrule
\endfoot
\emph{Quercus garryana} & MTWA & 16.4--23.8 & -0.6 / 8.0 & -0.2 / 7.6 & 0.0 / 7.5 \\
 & MTCO & 0.7--8.5 & 3.5 / 4.3 & 3.6 / 4.2 & 2.9 / 4.9 \\
 & GDD5 & 1503--3537 & 206 / 1828 & 273 / 1761 & 239 / 1795 \\
\emph{Arbutus menziesii} & MTWA & 16.0--22.3 & -0.2 / 6.5 & 0.2 / 6.1 & 0.3 / 5.9 \\
 & MTCO & 2.2--10.0 & 1.9 / 5.8 & 2.1 / 5.7 & 1.3 / 6.4 \\
 & GDD5 & 1604--3591 & 105 / 1882 & 171 / 1815 & 138 / 1849 \\
\emph{Rhododendron macrophyllum} & MTWA & 14.1--18.2 & 1.7 / 2.4 & 2.1 / 2.0 & 2.3 / 1.9 \\
 & MTCO & -1.9--8.8 & 6.0 / 4.6 & 6.2 / 4.5 & 5.4 / 5.2 \\
 & GDD5 & 1003--2643 & 706 / 934 & 773 / 867 & 739 / 901 \\
\end{longtable}
}


\figholder{figures/fig_oak_madrone.png}{Room to cool (blue, left) and to warm (orange, right) before the rainshadow sites leave the 5--95\% ranges of Garry oak and madrone, warmest month (left panel) and coldest month (right panel).}{fig:oakmadrone}

\figholder{figures/fig_envelopes.png}{Climate envelopes (5th--95th percentile) of the taxa that set or test the bounds, with site values (dashed).}{fig:envelopes}

\clearpage
\subsection{Does breadth track variability?}

As a test of the underlying proposition, the summer-temperature ranges of eastern lowland taxa are broader than those of western lowland and alpine taxa by every measure, including after allowing for range size (Table~\ref{tab:breadth}, Figure~\ref{fig:breadth}). The east--west difference in warmest-month breadth is 3.7\,\degC{} ($p=0.002$), 2.7\,\degC{} within latitude bands ($p=0.009$), and 2.0\,\degC{} after the range-size correction ($p=0.03$). Coldest-month breadth does not differ among groups. Western lowland and alpine taxa do not differ in summer breadth. With 36 tests, none passes a strict Bonferroni threshold of 0.0014, although the smallest $p$ (0.0017) is close and the tests are strongly correlated. The groups are small (6--9 taxa), and the eastern group is mostly shrubs while the others are mostly conifers, so plant form may account for part of the difference.

{\small
\begin{longtable}{llrrr}
\caption{Median climatic niche breadth (5th--95th percentile spread) by group. Local breadth: median within 1$^\circ$ latitude bands. Residual: breadth after regression on log range size \citep{slatyer2013}. Labrador tea and disturbance taxa excluded.}\label{tab:breadth}\\
\toprule
\textbf{Variable} & \textbf{Group} & \textbf{Breadth} & \textbf{Local breadth} & \textbf{Range-size residual} \\
\midrule
\endfirsthead
\toprule
\textbf{Variable} & \textbf{Group} & \textbf{Breadth} & \textbf{Local breadth} & \textbf{Range-size residual} \\
\midrule
\endhead
\bottomrule
\endfoot
MTWA & west lowland & 5.99 & 3.49 & -0.70 \\
 & east lowland & 9.72 & 6.28 & 1.26 \\
 & alpine ref. & 5.37 & 4.12 & -1.44 \\
MTCO & west lowland & 12.24 & 4.62 & 1.22 \\
 & east lowland & 15.69 & 5.48 & -0.29 \\
 & alpine ref. & 10.02 & 4.98 & -1.73 \\
MAT & west lowland & 7.52 & 3.37 & -0.42 \\
 & east lowland & 11.09 & 5.48 & 1.03 \\
 & alpine ref. & 6.66 & 4.74 & -1.29 \\
GDD5 & west lowland & 1671 & 705 & 113 \\
 & east lowland & 2588 & 1321 & 462 \\
 & alpine ref. & 913 & 679 & -783 \\
\end{longtable}
}


\figholder{figures/fig_breadth.png}{Warmest-month niche breadth by group, one point per taxon; bars are medians. Permutation $p$ for east vs.\ west shown in each panel.}{fig:breadth}

\clearpage
\subsection{Temperature and climate classes at the bounds}

Applying the regional scenarios to the WorldClim baseline gives maps of Washington at the coldest and warmest permitted Holocene temperatures (Figures~\ref{fig:mapmtwa} and~\ref{fig:mapkoppen}; Table~\ref{tab:koppen}). Only temperature is shifted; precipitation is held at its present values, so changes in class reflect temperature thresholds (mild-winter C versus cold-winter D versus alpine E, and the summer letters a, b, and c) and the temperature-dependent aridity threshold of the B class. If precipitation is unavailable, the maps fall back to temperature letters only. The scenario shift is derived from western sites and is extrapolated east of the Cascade crest.

\figholder{figures/fig_map_mtwa.png}{Warmest-month mean temperature in Washington at the coldest permitted Holocene bound, the present (1970--2000), and the warmest permitted bound. Hatched: east of the Cascade crest, extrapolated. Generated by the pipeline command \texttt{holocene-maps}.}{fig:mapmtwa}

\figholder{figures/fig_map_koppen.png}{\Kop{} classes for the same three scenarios, temperature shifted and precipitation held at present. Generated by \texttt{holocene-maps}.}{fig:mapkoppen}

{\small
\begin{longtable}{llr}
\caption{Share of area whose K\"{o}ppen--Geiger class (precipitation held at its present values) differs from the present under each temperature bound.}\label{tab:koppen}\\
\toprule
\textbf{Scenario} & \textbf{Region} & \textbf{\% area changing class} \\
\midrule
\endfirsthead
\toprule
\textbf{Scenario} & \textbf{Region} & \textbf{\% area changing class} \\
\midrule
\endhead
\bottomrule
\endfoot
Coldest Holocene & west of crest & 35.0 \\
Coldest Holocene & east of crest & 35.0 \\
Warmest Holocene & west of crest & 19.1 \\
Warmest Holocene & east of crest & 43.4 \\
\end{longtable}
}
{\small
\begin{longtable}{lrrr}
\caption{Statewide area (\%) in each K\"{o}ppen--Geiger class.}\label{tab:koppenarea}\\
\toprule
\textbf{Class} & \textbf{Coldest Holocene} & \textbf{Present} & \textbf{Warmest Holocene} \\
\midrule
\endfirsthead
\toprule
\textbf{Class} & \textbf{Coldest Holocene} & \textbf{Present} & \textbf{Warmest Holocene} \\
\midrule
\endhead
\bottomrule
\endfoot
BSk & 13.6 & 17.1 & 18.6 \\
BWk & 0.1 & 1.5 & 5.1 \\
Cfa & 0.0 & 0.2 & 0.6 \\
Cfb & 7.1 & 12.9 & 16.7 \\
Cfc & 1.2 & 0.5 & 0.3 \\
Csa & 0.0 & 1.2 & 5.7 \\
Csb & 14.2 & 17.6 & 21.9 \\
Csc & 0.2 & 0.0 & 0.0 \\
Dfa & 0.0 & 0.0 & 0.6 \\
Dfb & 9.5 & 14.2 & 13.1 \\
Dfc & 14.0 & 6.6 & 1.7 \\
Dsa & 0.0 & 0.5 & 2.6 \\
Dsb & 26.7 & 24.3 & 12.8 \\
Dsc & 9.7 & 3.3 & 0.3 \\
EF & 0.0 & 0.0 & 0.0 \\
ET & 3.6 & 0.1 & 0.0 \\
\end{longtable}
}


\section{Discussion}\label{sec:disc}

\paragraph{The answer so far.} The plant record bounds Holocene summer warmth in Washington's western lowlands at no more than about 2.5\,\degC{} above the 1970--2000 mean at the most constrained site, and 4--5\,\degC{} at the others, and bounds summer cooling at about 2.5\,\degC{} below it. The cold-side bound rests on a temperature-limited taxon and is the firmer of the two. The warm-side bound rests on moisture-limited conifers and should be read as provisional. The known warm periods, from the early-Holocene warm interval to the Medieval Climate Anomaly, lie within this range; the plants do not yet say where within it.

\paragraph{Consistency with other evidence.} The bounds agree in direction with the lake and glacier records: an early Holocene with expanded oak savanna, prairie, and fire in the lowlands \citep{whitlock1992,walsh2008battleground,brownhebda2002}, and smaller alpine glaciers than today \citep{menounos2009,zelenka2026olympic}, followed by cooling and the LIA glacier maximum. A cap of a few degrees on both sides is compatible with these changes, which were large in vegetation terms but did not remove the region's persistent conifers.

\paragraph{Why the cold side is better constrained.} The asymmetry has a simple cause. Subalpine taxa limited by summer warmth grow within a short distance of the lowland sites and never descended into them, which bounds cooling directly. No equivalent warm-adapted, temperature-limited taxon has yet been assessed as ``never present'' in the western lowlands. The Mediterranean-affinity trees that are present, Garry oak and madrone, have their warm edges far to the south and so constrain only the cold side.

\section{Next Steps: Firmer Warm-Side Bounds}\label{sec:next}

Three additions would turn the provisional warm cap into a firm one and might yield a lower bound.

First, \emph{warm-adapted taxa that never arrived}. Southern taxa whose northern limits lie in Oregon, and which could have migrated north during warmer intervals, would cap warming directly if they never reached Washington. Table~\ref{tab:candidates} lists candidates. Their use requires evidence that dispersal was not limiting over the Holocene and macrofossil, not pollen, evidence of absence, because several share pollen types with taxa already present.

{\small
\begin{longtable}{>{\raggedright\arraybackslash}p{0.24\textwidth}>{\raggedright\arraybackslash}p{0.22\textwidth}>{\raggedright\arraybackslash}p{0.20\textwidth}>{\raggedright\arraybackslash}p{0.24\textwidth}}
\caption{Candidate warm-side ``never present'' taxa for future analysis. Northern limits are approximate and must be recomputed from occurrences.}\label{tab:candidates}\\
\toprule
\textbf{Taxon} & \textbf{Approx.\ northern limit} & \textbf{Pollen identifiable?} & \textbf{Main caveat} \\
\midrule
\endfirsthead
\toprule
\textbf{Taxon} & \textbf{Approx.\ northern limit} & \textbf{Pollen identifiable?} & \textbf{Main caveat} \\
\midrule
\endhead
\bottomrule
\endfoot
\emph{Umbellularia californica} (California bay) & SW Oregon & poorly preserved & dispersal by birds and rodents \\
\emph{Calocedrus decurrens} (incense cedar) & N Oregon Cascades & no (Cupressaceae) & macrofossils required \\
\emph{Quercus kelloggii} (California black oak) & W Oregon & no (\emph{Quercus}) & lumped with Garry oak \\
\emph{Notholithocarpus densiflorus} (tanoak) & SW Oregon & partly & heavy acorns, slow migration \\
\emph{Pinus sabiniana}, \emph{P.\ attenuata} & N California / S Oregon & no (\emph{Pinus}) & edaphic specialists \\
\end{longtable}
}

Second, \emph{precipitation}. Adding monthly precipitation would show how much of the conifers' warm edges is a moisture edge, would allow full \Kop{} classes, and would let the eastern sites contribute, through the absence of ponderosa pine from Wildcat Lake.

Third, \emph{more sites and macrofossils}. Lost or gained taxa, which give lower bounds, require local presence and absence evidence such as macrofossils or small-hollow records \citep{sugimura2008,crausbay2017}. Verified coordinates for the Hoh and Wentworth records and a rerun on the 1~km climate grid are also needed.

\section{Limitations}\label{sec:limits}

The groups are small and the taxa were chosen, not sampled. The warm caps depend on one variable at one site, and on taxa limited mainly by moisture. Range envelopes assume rough equilibrium with modern climate \citep{svenningsandel2013}, and the realised range of a species can be narrower or wider than its tolerance. The bounds apply to sustained climate over many generations, not to individual decades. They are relative to 1970--2000; comparing them with later decades requires a multi-decade mean on the same baseline, and only an exceedance of the cap would be informative. Indigenous fire management shaped the prairies and oak savannas of the western lowlands for millennia \citep{leopoldboyd1999,pellattgedalof2014}, so their past extent is partly a cultural signal.

\section{Conclusion}

Plant range extremes offer a way to bound the warm side of Holocene climate, which glaciers record poorly. With 26 taxa and eight sediment records, the bounds for Washington's western lowlands are about $\pm2.5$\,\degC{} in summer relative to 1970--2000. The cold bound is firm, and the warm bound is provisional and set by a single site. No lower bound yet exists, because every documented change was one of abundance rather than presence. The method, the data, and the code are designed to be extended: more taxa, especially warm-adapted species that never reached Washington; precipitation; and local presence and absence records are the steps most likely to sharpen the warm-side answer.

\bibliographystyle{plainnat}
\bibliography{references}

\end{document}
